\section{Methods}
\label{sec:methods}

\subsection{Graph families and marked sets}
\label{sec:methods-graphs}

We use four synthetic graph families for the scaling experiments: complete
graphs $K_N$, cycles $C_N$, two-dimensional grids $L\times L$, and hypercubes
$Q_d$ with $N=2^d$. The multiple-target experiments add barbell graphs
(two cliques joined by a path) and $N\times N$ torus graphs with the marked-set
pattern used in the AGJK counterexample. The semantic experiment uses a
hand-built $19$-vertex concept graph with animal, vehicle, and instrument
communities, connected by the bridge concepts \texttt{bird} and
\texttt{airplane}. Decision experiments use a hub with one path per option;
matching depths of two branches can be coupled with a tunable weight.

For a graph $G=(V,E)$ we form the row-stochastic random-walk matrix
\begin{equation}
  P_{xy}=\frac{A_{xy}}{\sum_{z}A_{xz}},
\end{equation}
where $A$ is the adjacency matrix. All Szegedy and AGJK walks (E1, E2, E4,
E5) use the lazy walk $P\mapsto(1-\eta)P+\eta I$ with $\eta=0.1$; the AGJK
adversarial torus uses $\eta=0.2$, following the original construction.
Laziness removes the eigenvalue $-1$ of bipartite graphs and multiplies every
hitting time by exactly $1/(1-\eta)$, so it does not change any scaling law.
For the lazy walk we report the phase gap $1-\lambda_2$ of the discriminant
matrix, because on bipartite graphs the absolute gap
$1-\max_{i\ge2}\lvert\lambda_i\rvert$ is capped at $2\eta$ by the image
$2\eta-1$ of the eigenvalue $-1$. (When $\lambda_2<0$, as on small complete
graphs, the phase gap exceeds $1$.) A
marked set $M\subset V$ is either a single target, an independent set of $k$
vertices chosen by farthest-point sampling, a connected cluster grown by BFS,
or a set of option terminals. Table~\ref{tab:experiment-matrix} lists the
experiment matrix.

\begin{table}[t]
  \centering
  \small
  \begin{tabular}{@{}>{\raggedright\arraybackslash}p{0.05\linewidth}
    >{\raggedright\arraybackslash}p{0.23\linewidth}
    >{\raggedright\arraybackslash}p{0.24\linewidth}
    >{\raggedright\arraybackslash}p{0.15\linewidth}
    >{\raggedright\arraybackslash}p{0.25\linewidth}@{}}
    \toprule
    ID & Graph / system & Marked set & Noise & Main metric \\
    \midrule
    E1 & $K_N$, $C_N$, grid, $Q_d$ & single & none & scaling of $T_q$ and $\HT$ \\
    E2 & grid, barbell & single / independent / clustered & none & $p_M(t)$, $\HT^+$ \\
    E2b & adversarial torus & AGJK pattern & none & $\HT^+/\HT$ \\
    E3 & toy semantic graph & single target & none & activation, AUC \\
    E4 & multi-branch graph & option terminals & none & choice prob., RT \\
    E5 & grid, toy semantic & single, clustered & partial measurement & timed peak, detection by $T$ \\
    \bottomrule
  \end{tabular}
  \caption{Experiment matrix. Each row corresponds to a processed result file in
  \texttt{code/results/processed/}.}
  \label{tab:experiment-matrix}
\end{table}

\subsection{Classical hitting time}
\label{sec:methods-classical}

For a marked set $M$, let $h_x$ be the expected time to reach $M$ from
$x\notin M$, with $h_x=0$ for $x\in M$. The hitting times solve the linear
system
\begin{equation}
  h_x=1+\sum_{y}P_{xy}h_y,\qquad x\notin M.
  \label{eq:hitting-linear}
\end{equation}
The classical hitting time is the stationary average conditioned on starting
outside $M$,
\begin{equation}
  \HT(P,M)=\frac{\sum_{x\notin M}\pi_x h_x}{1-p_M},
  \qquad p_M=\sum_{x\in M}\pi_x.
  \label{eq:hitting}
\end{equation}
We solve \eqref{eq:hitting-linear} directly for all grids and barbell graphs;
the synthetic families are small enough that dense linear algebra is
inexpensive.

\subsection{Szegedy and AGJK search protocol}
\label{sec:methods-szegedy}

The quantum walk uses the register convention of AGJK,
\begin{equation}
  V(P)\lvert\bar 0\rangle\lvert x\rangle
  =\sum_y\sqrt{P_{xy}}\lvert y,x\rangle,\qquad
  W(P)=V^\dagger(P)\,\mathrm{Shift}\,V(P)\,\mathrm{Ref},
\end{equation}
with $\mathrm{Shift}\lvert x,y\rangle=\lvert y,x\rangle$ and
$\mathrm{Ref}=(2\lvert\bar 0\rangle\langle\bar 0\rvert-I)\otimes I$. The
discriminant matrix is $D_{xy}=\sqrt{P_{xy}P_{yx}}$.

For multiple or clustered targets we use the interpolated walk
\begin{equation}
  P(s)=(1-s)P+sP',\qquad
  P'=\begin{pmatrix}P_{UU}&P_{UM}\\0&I\end{pmatrix},
  \label{eq:interpolated}
\end{equation}
and the corresponding walk operator $W(s)$ with an ancilla that records whether
the current vertex is marked. The initial state is
$\lvert\bar 0\rangle\lvert\sqrt{\pi}\rangle$. For each graph and marked set we
scan $s$ over a fixed grid and select the value that maximizes the AGJK lower
bound
\begin{equation}
  q_t(s)=\left\lVert\Pi_M T_t(D(s))\lvert\sqrt{\pi}\rangle\right\rVert^2,
  \qquad t\le 3\sqrt{\HT}+5,
  \label{eq:lower-bound}
\end{equation}
where $T_t$ is the Chebyshev polynomial of the first kind. We then run the
actual walk $W(s)$ and record
\begin{equation}
  p_M(t)=\left\lVert(I\otimes\Pi_M)W^t(s)\lvert\bar 0\rangle
  \lvert\sqrt{\pi}\rangle\right\rVert^2.
\end{equation}
The quantum search time is $T_q=\arg\max_{t\ge1}p_M(t)$ and the peak success
probability is $p_{\max}=\max_t p_M(t)$. The extended hitting time $\HT^+$ is
computed from the spectral definition
\begin{equation}
  \HT(s)=\frac{1}{1-p_M}\sum_{k=1}^{n-1}
  \frac{\lvert\langle v_k(s)\rvert\sqrt{\pi_U}\rangle\rvert^2}{1-\lambda_k(s)},
  \qquad
  \HT^+=\lim_{s\to1}\HT(s),
  \label{eq:ht-plus}
\end{equation}
evaluated at $s=1-10^{-6}$ as a numerical proxy. Equation
\eqref{eq:ht-plus} is used only for diagnostics; the search curves use the
actual operator $W(s)$.

\subsection{Semantic activation protocol}
\label{sec:methods-semantic}

For the semantic experiment we use a continuous-time quantum walk (CTQW) and a
continuous-time random walk (CRW) from the same cue vertex. Let
\begin{equation}
  L_{\mathrm{sym}}=I-D^{-1/2}AD^{-1/2}
\end{equation}
be the symmetric normalized Laplacian and
$L_{\mathrm{rw}}=I-D^{-1}A$ the random-walk generator. The two models are
\begin{equation}
  \lvert\psi(t)\rangle=e^{-iL_{\mathrm{sym}}t}\lvert c\rangle,
  \qquad
  \rho(t)=\rho(0)e^{-L_{\mathrm{rw}}t},
  \label{eq:ctqw-crw}
\end{equation}
with $\rho(0)=\lvert c\rangle\langle c\rvert$ at the cue $c$. Time is reported
in units of the spectral gap, $\tau=\lambda_1t$, so that both walks are
compared at the same relaxation scale. The target activation is
$p_T(t)=\lvert\langle T\vert\psi(t)\rangle\rvert^2$ for CTQW and
$p_T(t)=\rho_T(t)$ for CRW. We record the peak activation, time of peak, area
under the activation curve, entropy at the peak, half-peak coverage, and the
expected graph distance from the cue at the peak.

\subsection{Decision protocol}
\label{sec:methods-decision}

The decision graph has a central hub and one path per option. Each option
$i$ corresponds to a terminal set $M_i=\{\text{terminal}_i\}$. To control
similarity, matching depths of branches $i$ and $j$ are joined by an edge of
weight $c$; $c=0$ gives independent branches and large $c$ makes the two
options highly similar. The quantum model runs the interpolated walk $W(s)$
with $s=0.9$ and initial state $\lvert\bar 0\rangle\lvert\text{hub}\rangle$.
The classical model runs the same graph with all option terminals absorbing.
For each model we compute the option probabilities $p_i(t)$ and the total
decision probability $P_M(t)=\sum_i p_i(t)$, with $T_{\max}=400$. We compare
four readouts. (i) \emph{First crossing} (our original readout): choice
probabilities are the normalized $p_i(t)$ at the first $t$ with
$P_M(t)\ge0.5$, or at $\arg\max_tP_M(t)$ if the threshold is never reached.
(ii) \emph{Random stopping}: $P_i\propto T_{\max}^{-1}\sum_{t=1}^{T_{\max}}p_i(t)$.
(iii) \emph{Detection}: the measurement model of \cref{sec:methods-noise} with
rate $p=0.2$, where a measurement landing on terminal $i$ ends the decision
with choice $i$; $P_i$ is the probability of choice $i$ by $T_{\max}$
(essentially one) and the decision time is the mean first-detection time.
(iv) \emph{Classical}: absorption probabilities and mean absorption time of the
classical walk. To test robustness we repeat every configuration for
$\eta\in\{0,0.05,0.1,0.2\}$ and for ten independent draws in which every edge
weight is multiplied by a factor uniform in $[0.8,1.2]$, and report the spread
(max minus min) of the probability of the uncoupled option(s).

\subsection{Partial-measurement protocol}
\label{sec:methods-noise}

The decoherence experiment uses the Szegedy walk $W(0)$ of the lazy chain $P_\eta$ ($\eta=0.1$), i.e.
interpolation $s=0$, from a cue
vertex. Before each step, with probability $p$ we measure the vertex register
in the position basis and reset the state to
$\lvert\bar 0\rangle\lvert x\rangle$; with probability $1-p$ the state evolves
unitarily. Thus $p=0$ is the pure quantum walk and $p=1$ is exactly the
classical random walk on the same graph. We compute the exact average over
measurement trajectories by precomputing the position marginals and marked
probabilities of the basis-state ensembles
$W(0)^a\lvert\bar 0\rangle\lvert x\rangle$, then propagating the age
distribution of the most recent measurement. This avoids Monte Carlo noise and
makes the sweep deterministic. The success probability is the probability that
the walker is at a marked vertex after the step. We use $T_{\max}=120$ and
$p\in\{0,0.05,0.1,0.2,0.3,0.5,0.75,1\}$.

Because $p_M(t)$ oscillates at small $p$, we do not compare rates at a single
deadline. In the \emph{readout model} we report the peak, the time-averaged
success $\bar p(H)=H^{-1}\sum_{t=1}^{H}p_M(t)$ (the success under a uniform
deadline on $\{1,\dots,H\}$, the same device as randomized stopping in quantum
search with unknown optimal time), and the coefficient of variation of $p_M$.
In the \emph{detection model}, a measurement that lands on $M$ ends the search;
the recursion is the same except that this mass is moved to a ``found''
counter instead of being reset (\cref{app:measurement}). With a final readout
at the deadline $T$, $F_p(T)$ is the probability that the target has been found
by $T$; $F_0$ equals the pure readout curve and $F_1(T)=\Pr[\tau_M\le T]$ for
the classical walk. We report $F_p(T)$ for $T\in\{20,40,80,120\}$ and the best
rate $p^\ast(T)=\arg\max_pF_p(T)$.

\subsection{Metrics and reproducibility}
\label{sec:methods-reproducibility}

The main metrics are the classical hitting time $\HT$, the extended hitting
time $\HT^+$, the quantum search time $T_q$, the peak success probability
$p_{\max}$, the speedup $\HT/T_q$, the target AUC, the activation entropy and
radius, choice probabilities, decision time, and the success probability at a
fixed deadline. Scaling exponents are fitted by ordinary least squares in
log--log coordinates.

The implementation is in Python~3.12 with NumPy, SciPy, NetworkX, Matplotlib,
PyYAML, and pytest. The environment is created with the Miniforge file
\texttt{code/environment.yml} (environment \texttt{qwalk-cognition}). All
random choices use explicit seeds; the configuration files are
\texttt{code/configs/base.yaml} and \texttt{code/configs/experiments.yaml}. The
test suite contains 61 tests covering unitarity, norm preservation, the
classical limit, hitting-time identities, decision normalization, and the
exact measurement ensemble. Every figure in this paper is regenerated by the
commands listed in \texttt{code/README.md}.

\FloatBarrier
